%! Author = sriadityavedantam
%! Date = 3/30/26

\section{Polynomials}

\begin{definition}[Polynomial]
    A polynomial with coefficients in $\ri$ is denoted by $\ri[x]$.
    It consists of expressions of the form
    \[
        a_0 + a_{1}x + a_{2}x^2 + \hdots + a_{n}x^n,
    \]
    where $a_i\nri$.
    We can think of the $a_i$ as coefficients.
\end{definition}

\begin{proposition}{
    Addition and multiplication in $\ri[x]$ are defined component-wise.
}
    Let
    \begin{align*}
        f(x) &= a_0 + \hdots + a_{n}x^n,\\
        g(x) &= b_0 + \hdots + b_{m}x^m,
    \end{align*}
    where $m\geq n$.
    Then
    \[
        f(x)+g(x) = (a_0+b_0) + \hdots + (a_n+b_n)x^n + b_{n+1}x^{n+1} + \hdots + b_{m}x^m.
    \]
    Multiplication is defined by distributing and collecting like terms.
\end{proposition}

This informal definition raises several questions.
What is $x$?
Is it an element of $\ri$?
If not, what does it mean to multiply $x$ by a ring element?
To answer these questions, note that an expression of the form
\[
    a_0 + a_{1}x + a_{2}x^2 + \ldots + a_{n}x^n
\]
makes sense provided that the $a_i$ and $x$ are all elements of some larger ring.
An analogy might be helpful here.
The number $\pi$ is not in the ring of integers $\mathbb{Z}$, but expressions such as
\[
    3 - 4\pi + 12\pi^2 + \pi^3
\]
make sense in $\mathbb{R}$.
Furthermore, it is not difficult to verify that the set of all numbers of the form
\[
    \sum_{i=0}^n a_i\pi^i,
\]
with $n\geq 0$ and $a_i\in\mathbb{Z}$, is a subring of $\mathbb{R}$ that contains both $\mathbb{Z}$ and $\pi$.
For the present, we shall think of polynomials with coefficients in a ring $\ri$ in much the same way, as elements of a larger ring that contains both $\ri$ and a special element $x$ that is not in $\ri$.
Feel free to check that $\ri[x]$ is a ring, but we will be concentrating on $\z[x],\Q[x],\r[x],\c[x],\zp[x]$, and have their elements denoted by $f(x)$ or $P(x)$.

\begin{definition}[Degree of a Polynomial]
    If $f(x)\in\ri[x]$, the degree of $f(x)$, denoted by $\deg f(x)$, is the largest $n$ for which the coefficient of $x^n$ is not $0$.
    The coefficient $a_n$ is called the leading coefficient.
\end{definition}

\begin{definition}[Additive Identity of $\ri[x]$]
    The additive identity in $\ri[x]$ is the zero polynomial, where all coefficients are $0$.
\end{definition}

If $\deg f(x)=0$, then $f(x)$ is a constant polynomial.

\begin{proposition}[Degree Arithmetic]{
    Suppose $\deg f(x)=m$ and $\deg g(x)=n$.
}
    \[
        \deg(f(x)+g(x)) \leq \max\{m,n\}.
    \]
    If $m\neq n$, then
    \[
        \deg(f(x)+g(x)) = \max\{m,n\}.
    \]
    If $m=n$, then
    \[
        \deg(f(x)+g(x)) \leq m,
    \]
    with equality unless cancellation occurs.
    Also,
    \[
        \deg(f(x)g(x)) \leq \deg f(x) + \deg g(x).
    \]
    If $\ri$ is an integral domain, then
    \[
        \deg(f(x)g(x)) = \deg f(x) + \deg g(x).
    \]
\end{proposition}

Let $f(x),g(x)\in\z_4[x]$.
Let $f(x)=2x$ and $g(x)=2x^2$.
Then
\[
    f(x)g(x)=4x^3=0.
\]

    {\Huge{From now on $\ri$ is an Integral Domain.}}\\

Given that $\ri$ is an integral domain, one may naturally ask what are the units of $\ri[x]$.

\begin{lemma*}{
    Suppose $u(x)$ is a unit with multiplicative inverse $v(x)$.
    Then
    \[
        u(x)v(x)=1=1+0x+0x^2+\hdots.
    \]
}
\end{lemma*}